\chapter{Краткий справочник команд}

\section{Версия}

\begin{lstlisting}
python -c "import grandgousier; print(grandgousier.__version__)"
\end{lstlisting}

\section{Проверка Python-кода}

\begin{lstlisting}
python -m compileall grandgousier
pytest -q
\end{lstlisting}


\section{Canonical Stage 1 plan (0.6.0)}

Metadata-only план без скачивания science pixels и без записи в БД:

\begin{lstlisting}
python -m grandgousier.cli stage1-plan `
  --config grandgousier.yaml `
  --source robophot-local `
  --filter i `
  --out reports\stage1_plan.json `
  --csv-out reports\stage1_plan.csv
\end{lstlisting}

Source задаётся в секции \texttt{archives} application config, а derived products --- отдельно в \texttt{products}.

\section{Canonical Stage 1 execute (0.6.1.x)}

Явная materialization/reduction в ProductStore без автоматического Stage~2:

\begin{lstlisting}
python -m grandgousier.cli stage1-execute `
  --config grandgousier.yaml `
  --source master-prepared `
  --limit 1 `
  --plan-out reports\stage1_execute_plan.json `
  --report-out reports\stage1_execution.json `
  --csv-out reports\stage1_execution.csv
\end{lstlisting}

\cli{--overwrite} разрешает заменить существующий product только явным действием оператора. Без него существующий canonical product переиспользуется лишь при совпадении fingerprint и checksum. Legacy MASTER output без sidecar не принимается автоматически: его path-valued provenance cards не доказывают неизменность байтов raw/header/dark/flat, поэтому требуется явная пересборка с \cli{--overwrite}. Byte-identical processed import можно безопасно принять сравнением checksum. Plan с любым \cli{BLOCKED} item не выполняется частично.

\section{Stage 1 -> Stage 2 exact handoff (0.6.2.x)}

Read-only preflight по exact products одного execution report:

\begin{lstlisting}
python -m grandgousier.cli stage1-handoff `
  --config grandgousier.yaml `
  --execution-report reports\stage1_execution.json `
  --report-out reports\stage1_handoff.json
\end{lstlisting}

После проверки:

\begin{lstlisting}
python -m grandgousier.cli stage1-handoff `
  --config grandgousier.yaml `
  --execution-report reports\stage1_execution.json `
  --execute `
  --stage2-manifest-out reports\stage1_handoff_stage2.json `
  --report-out reports\stage1_handoff_execute.json
\end{lstlisting}

Команда не сканирует ProductStore. Она валидирует product/provenance checksum, выполняет P4 preflight до индексации и передаёт Stage~2 только exact paths со статусом \texttt{SELECTED\_FOR\_STAGE2}. \texttt{ALREADY\_REPRESENTED} не создаёт новый Frame и не запускает Stage~2. В 0.6.2.1 исторический \texttt{EXISTING\_MULTIPLE\_REPRESENTATIVES} также сохраняется как handoff no-op: diagnostic P4 action остаётся в отчёте, существующие representatives не меняются, новый Frame не создаётся.

\section{Stage 1A: найти наблюдения в MASTER}

\begin{lstlisting}
python -m grandgousier.cli master-archive-search `
  --search-transport http `
  --http-user "googlov" `
  --ra 327.60625 `
  --dec 72.1686111 `
  --radius-deg 0.8 `
  --out "C:\master\gpx\inventory.csv"
\end{lstlisting}

HTTPS discovery использует штатный spatial search MASTER; для \cli{--search-transport http} сейчас требуется \cli{--radius-deg <= 0.95}.

\section{Stage 1A: скачать raw/header/dark/flat}

Plan-only:

\begin{lstlisting}
python -m grandgousier.cli master-archive-fetch `
  --manifest "C:\master\gpx\inventory.csv" `
  --out-root "C:\master\gpx"
\end{lstlisting}

Execution:

\begin{lstlisting}
python -m grandgousier.cli master-archive-fetch `
  --manifest "C:\master\gpx\inventory.csv" `
  --out-root "C:\master\gpx" `
  --http-user "MASTER_USER" `
  --execute
\end{lstlisting}

Для HTTP transfer credentials задаются только через environment \cli{GRANDGOUSIER\_MASTER\_HTTP\_USER} и \cli{GRANDGOUSIER\_MASTER\_HTTP\_PASSWORD}; password не передаётся аргументом CLI. Downloader проверяет \cli{Content-Length}/\cli{Content-Range}, итоговый размер и resume offset. При сетевой ошибке или stall сохраняется \cli{.part}, а повтор использует Range от фактического offset. По умолчанию stall timeout равен 300 секундам, progress interval --- 15 секундам; параметры можно переопределить в конфигурации archive assets.


\section{Показать provisional серии MASTER до скачивания}

\begin{verbatim}
python -m grandgousier.cli master-archive-series `
  --manifest "C:\master\gpx\full_inventory.csv" `
  --max-field-separation-arcmin 30 `
  --max-series-gap-minutes 60
\end{verbatim}

\section{Stage 1B: построить processed MASTER FITS}

Plan-only:

\begin{lstlisting}
python -m grandgousier.cli master-archive-reduce `
  --manifest "C:\master\gpx\inventory.csv" `
  --cache-root "C:\master\gpx"
\end{lstlisting}

Execution:

\begin{lstlisting}
python -m grandgousier.cli master-archive-reduce `
  --manifest "C:\master\gpx\inventory.csv" `
  --cache-root "C:\master\gpx" `
  --execute
\end{lstlisting}

Для smoke test можно добавить \cli{--limit 1}. По умолчанию выходной root --- \cli{<cache-root>/processed}; существующие processed FITS не перезаписываются без \cli{--overwrite}.

Компактировать legacy uint16 DQ и удалять raw после подтверждённого Stage~1B:

\begin{lstlisting}
python -m grandgousier.cli master-archive-reduce `
  --manifest "C:\master\gpx\inventory.csv" `
  --cache-root "C:\master\gpx" `
  --compact-existing `
  --delete-raw-after-success `
  --execute
\end{lstlisting}

Безопасная очистка уже обработанных файлов сначала запускается plan-only:

\begin{lstlisting}
python -m grandgousier.cli master-archive-prune `
  --manifest "C:\master\gpx\inventory.csv" `
  --cache-root "C:\master\gpx" `
  --scope raw

python -m grandgousier.cli master-archive-prune `
  --manifest "C:\master\gpx\inventory.csv" `
  --cache-root "C:\master\gpx" `
  --scope processed
\end{lstlisting}

Фактическое удаление требует одновременно \cli{--execute --yes}.

\section{Планирование series processing}

\begin{lstlisting}
python -m grandgousier.cli process-archive-path --path "D:\archive\series"
\end{lstlisting}

\section{Обработка series}

\begin{lstlisting}
python -m grandgousier.cli process-archive-path `
  --path "D:\archive\series" `
  --execute `
  --manifest-out reports\series.json
\end{lstlisting}

\subsection{Selective rerun и conventional ensemble threshold (0.5.45.1)}

Для повторной обработки только существующей series среди кадров указанного пути:
\begin{lstlisting}
python -m grandgousier.cli process-archive-path `
  --path "D:\archive\night" `
  --execute `
  --series-id 78 `
  --ensemble-min-ref-stars-per-frame 20
\end{lstlisting}
\texttt{--series-id} допускается несколько раз. В этом режиме новые field series не строятся.
Default \texttt{--ensemble-min-ref-stars-per-frame} равен 30 и не был изменён в 0.5.45.1.
Если conventional ensemble не создаёт usable frame offsets, pipeline завершает run как
\path{FAILED_PRECONDITION_ENSEMBLE} и сохраняет base-ensemble telemetry.

\section{Сборка переносного developer context}

\begin{lstlisting}
python scripts\build_context_bundle.py
\end{lstlisting}

Эта команда относится к сопровождению проекта, а не к научной обработке данных, но полезна при переносе разработки в новый чат.

\section{Stage 3: проверить готовность поля к межсерийной калибровке}

\begin{lstlisting}
python -m grandgousier.cli plan-interseries-calibration `
  --field-id 12 `
  --min-points-per-object 5 `
  --min-common-objects 20 `
  --required-stage2-recipe stage3_scale_v1
\end{lstlisting}

Команда read-only: она выбирает опубликованные SCIENCE generations, считает common objects и строит граф совместимых series, но ещё не вычисляет и не записывает photometric offsets. Опциональный \texttt{--required-stage2-recipe} дополнительно проверяет persisted producing-run provenance и исключает \texttt{RECIPE\_MISMATCH}/\texttt{MISSING\_PROVENANCE} из READY-графа.

\section{SCIENCE storage в clean schema 0.7}

Legacy maintenance commands для \texttt{science\_measurements}, rebuildable
\texttt{science\_photometry\_cache} и historical flag overrides в 0.7 удалены.
Новые SCIENCE products штатно записываются в durable
\texttt{science\_photometry}; отдельный backfill/retirement workflow для новой
БД не требуется. Frozen 0.6 reference database не должна обслуживаться CLI 0.7.

\section{Stage 3: read-only M1--M3 межсерийная калибровка}

\begin{lstlisting}
python -m grandgousier.cli analyze-interseries-calibration `
  --field-id 12 `
  --filter V `
  --min-points-per-object 5 `
  --min-common-objects 20 `
  --out-root reports\stage3_field12_V
\end{lstlisting}

Для неоднозначного поля можно дополнительно указать \texttt{--instrument-unit-id}, \texttt{--catalog-name} или \texttt{--reference-series-id}. Команда read-only и не создаёт Stage~3 rows в БД.

\section{Stage 3 M4: persist immutable graph generation}

\begin{lstlisting}
python -m grandgousier.cli persist-interseries-calibration `
  --field-id 33 `
  --filter R `
  --catalog-name "auto_catalog_master_20180907_GPX_TF17_2_mixed" `
  --reference-series-id 77 `
  --min-points-per-object 5 `
  --min-common-objects 20
\end{lstlisting}

Перед первым запуском требуется \texttt{migrate-db}. Повтор той же команды idempotent. Для connected tree без closure требуется явный \texttt{--allow-sparse-graph}.

\section{Stage 3 M4: show/verify persisted generation}

\begin{lstlisting}
python -m grandgousier.cli show-interseries-calibration --field-id 33
python -m grandgousier.cli show-interseries-calibration --generation-id 1
python -m grandgousier.cli verify-interseries-calibration --generation-id 1
\end{lstlisting}

\section{Stage 3 M5: export global object light curve + plots}

\begin{lstlisting}
python -m grandgousier.cli export-global-object-light-curve `
  --object-id 12345 `
  --field-id 33 `
  --generation-id 1 `
  --generation-id 2 `
  --out-root reports\object_12345_global
\end{lstlisting}

По умолчанию основной CSV содержит exact accepted Stage~2 SCIENCE rows на M4 global scale. Plot строится отдельно для каждого filter и консервативно переключается на биннинг, когда exact linked source-row count до QC превышает 5000. Полезные опции:

\begin{lstlisting}
--plot-max-points 5000
--plot-reduction auto|bin|none
--plot-bin-minutes 5
--plot-time-axis auto|continuous|compressed
--plot-format png|pdf|svg
--include-rejected
--include-rejected-debug
--allow-partial
--no-plots
\end{lstlisting}

\texttt{--allow-partial} не применяет offset к несовместимым series: он только разрешает исключить их из export с явным warning. \texttt{--include-rejected-debug} пишет отдельный diagnostic CSV. При exact-linked source rows, но нуле accepted rows, M5 0.5.47.1 возвращает status \texttt{NO\_ACCEPTED\_POINTS\_AFTER\_QC} и exit code 2. Начиная с 0.5.47.2 \texttt{--plot-max-points} является hard cap на filter, а \texttt{--plot-time-axis auto} по умолчанию визуально сокращает большие пустые gaps между series, не изменяя MJD в CSV. При числе series больше 12 включается archive-overview без per-series legend.

Начиная с 0.5.47.3 default M5 saturation veto является epoch-level: \texttt{SATURATED\_OR\_NEAR\_SATURATED} отвергает конкретную measurement, а object-level \texttt{SATURATED\_LIKE} остаётся diagnostic-only для M5. Общий Stage~2/Object-QC policy не меняется. При необходимости strict series-level rejection можно вернуть явно: \texttt{--extra-reject-flags SATURATED\_LIKE}.


\section{M8: resumable processing всего processed-архива поля}

\begin{lstlisting}
python -m grandgousier.cli process-field-archive `
  --processed-root "C:\master\gpx\smoke\processed" `
  --field-id 33 `
  --catalog-name "auto_catalog_master_20180907_GPX_TF17_2_mixed" `
  --stage3-policy "V:15:88" `
  --stage3-policy "R:20:89" `
  --control-object-id 30112 `
  --control-object-id 29786 `
  --resume `
  --execute `
  --out-root reports\gpx_tf17_final
\end{lstlisting}

Начиная с 0.5.51.0 эта команда проходит через service layer. Дополнительно поддерживаются \texttt{--job-config FILE}, \texttt{--plan-only} и diagnostic-only \texttt{--no-preflight}. До M8 выполняются readiness checks и PostgreSQL advisory lock для field/catalog/root scope. Machine-readable job metadata сохраняется в \texttt{<out-root>/\_jobs/<job-id>}.

Отдельная deployment-проверка:

\begin{lstlisting}
python -m grandgousier.cli field-archive-readiness `
  --processed-root "C:\master\gpx\smoke\processed" `
  --instrument-profile MASTER `
  --out-root reports\gpx_tf17_final
\end{lstlisting}

Без \texttt{--execute} выполняется plan-only. Policy задаётся явно и повторяется для каждого filter/system; M8 v1 требует одну policy на filter. Базовая форма --- \texttt{FILTER:MIN\_COMMON}; после неё можно добавить reference series ID и instrument unit ID. \texttt{--include-date} задаёт allow-list, \texttt{--exclude-date} --- deny-list; оба параметра повторяемые. \texttt{--no-resume} запрещает повторное использование уже индексированных/incomplete batches. Начиная с 0.5.54.3 resume readiness для P4-linked данных проверяет representation physical Observation в READY series целевого field, поэтому alternate Frame products не вызывают ложный \texttt{RESUME\_INCOMPLETE}; genuinely unrepresented Observation и legacy NULL identity остаются blocking. Leaf batch с несколькими календарными датами останавливается как conflict. Начиная с 0.5.48.1 заведомо короткий новый/incomplete batch, где максимальная series-like группа содержит менее 10 FITS, помечается \texttt{SKIP\_TOO\_SMALL} и не передаётся в Stage~2.

\section{Audit и очистка derived Stage 2 данных}

Dry-run:

\begin{lstlisting}
python -m grandgousier.cli reset-derived-data `
  --scope stage2 `
  --all-series
\end{lstlisting}

Destructive execution:

\begin{lstlisting}
python -m grandgousier.cli reset-derived-data `
  --scope stage2 `
  --all-series `
  --execute `
  --yes
\end{lstlisting}

\section{Stage 3: audit catalog identity}

\begin{lstlisting}
python -m grandgousier.cli audit-field-catalog-identity `
  --field-id 2
\end{lstlisting}

Для первого same-system Stage~3 MVP повторные series должны иметь общий catalog namespace.


\paragraph{Экспериментальная политика aperture correction.}
Для \texttt{process-archive-path} и \texttt{run-field-pipeline} доступен диагностический параметр:
\begin{lstlisting}
--aperture-correction-plane-policy evidence
--aperture-correction-plane-policy safety_only
\end{lstlisting}
Значение \texttt{evidence} является стандартным; \texttt{safety\_only} предназначено только для контролируемого A/B-теста пространственной модели.

\section{PSF field diagnostics}

Диагностика уже существующей серии без повторной фотометрии:
\begin{lstlisting}
python -m grandgousier.cli psf-field-diagnostics `
  --series-id 77 `
  --out-dir "exports_psf\series77"
\end{lstlisting}

Основные параметры отбора:
\begin{lstlisting}
--max-sources-per-frame 500
--min-snr 20
--cutout-radius-pix 7
--min-isolation-pix 12
--clip-sigma 4
--n-iterations 3
\end{lstlisting}

В обычном \texttt{process-archive-path} диагностический шаг включается так:
\begin{lstlisting}
--enable-psf-field-diagnostics
\end{lstlisting}
Для этой команды PSF diagnostics автоматически требует source detection, но не включает
обязательное detection-to-catalog matching, если отдельно не указан
\texttt{--enable-detection-branch}. Science photometry и calibration semantics от включения
PSF diagnostics в версии 0.5.44.0 не меняются.

\subsection{adaptive-aperture-diagnostics}
Экспериментальная A/B фотометрия circular versus local-PSF adaptive ellipse для существующей
field series. Команда не изменяет science measurements и не пишет новый calibration product.

\begin{verbatim}
python -m grandgousier.cli adaptive-aperture-diagnostics `
  --series-id 77 `
  --out-dir exports_adaptive\series77_R
\end{verbatim}

Наиболее важные параметры: \texttt{--base-aperture-radius-pix},
\texttt{--aperture-fwhm-factor}, \texttt{--interpolation-neighbors},
\texttt{--max-axis-ratio}, \texttt{--max-radius-scale}. Для синхронного A/B сравнения MASTER
следует отдельно запускать EAST/V и WEST/R series и сравнивать object-level RMS/MAD versus
normalized detector Y.

\section{Выбор базового flux estimator Stage 2}

Для \texttt{process-archive-path}, \texttt{run-field-pipeline} и прямого \texttt{run-forced-photometry}:
\begin{lstlisting}
--photometry-estimator circular
--photometry-estimator adaptive_ellipse
--photometry-estimator spatial_psf
--photometry-estimator empirical_psf
\end{lstlisting}
В source-кандидате 0.7.0.10 для новых запусков Stage~2 без этого параметра default ---
\texttt{adaptive\_ellipse} для каждого источника (с явно отмеченным circular fallback при
отсутствии надёжной локальной геометрии). Для прежнего circular-поведения укажите
\texttt{--photometry-estimator circular} явно в командах выше; для
\texttt{process-archive-path} без \texttt{--execute} остаётся только план.
Начиная с 0.5.44.3 summary команды \texttt{process-archive-path} дополнительно показывает effective
\texttt{photometry\_estimator} для каждой pipeline-серии. Режим
\texttt{adaptive\_ellipse} автоматически требует source detection и строит локальную PSF-модель
по bright isolated detections. Основные параметры настройки:
\begin{lstlisting}
--adaptive-min-psf-snr 20
--adaptive-max-psf-sources-per-frame 500
--adaptive-psf-cutout-radius-pix 7
--adaptive-psf-min-isolation-pix 12
--adaptive-interpolation-neighbors 24
--adaptive-min-psf-references 20
--adaptive-aperture-fwhm-factor 1.0
--adaptive-max-axis-ratio 3.0
--adaptive-max-radius-scale 2.5
--no-adaptive-circular-fallback
\end{lstlisting}
Последний параметр предназначен прежде всего для строгих A/B-тестов: без него при невозможности
построить adaptive geometry используется явно отмеченная circular fallback photometry.

\subsection{Spatial PSF estimator}
Для полного Stage~2 пространственно меняющаяся profile-fit фотометрия включается так:
\begin{lstlisting}
--photometry-estimator spatial_psf
\end{lstlisting}
Основные параметры:
\begin{lstlisting}
--psf-profile moffat
--psf-moffat-beta 2.5
--psf-fit-radius-fwhm 2.5
--psf-min-fit-radius-pix 5
--psf-max-fit-radius-pix 14
--psf-centroid-search-radius-pix 0.75
--psf-centroid-grid-steps 3
--psf-residual-clip-sigma 5
--psf-residual-clip-iterations 2
--no-psf-adaptive-fallback
\end{lstlisting}
По умолчанию при локальной неудаче PSF fit используется adaptive ellipse; при отсутствии local
PSF field возможен явно помеченный circular fallback. Для строгого A/B оба fallback можно
отключить соответствующими параметрами.


\subsection{Empirical spatial PSF estimator}
Для A/B без предположения аналитической формы PSF:
\begin{lstlisting}
--photometry-estimator empirical_psf
--empirical-psf-grid-size 3
--empirical-psf-stamp-radius-pix 8
--empirical-psf-reference-radius-pix 5.0
--empirical-psf-references-per-node 48
--empirical-psf-min-references-per-node 12
--no-empirical-psf-adaptive-fallback
\end{lstlisting}
Шаблоны строятся из зарегистрированных реальных stellar stamps и нормируются к указанной
reference aperture. По умолчанию локальная неудача empirical fit переходит в adaptive ellipse,
а затем при необходимости в явно отмеченный circular fallback.


\subsection{Spatial ensemble calibration}
Для experimental residual field correction после обычного ensemble:
\begin{lstlisting}
--enable-spatial-ensemble
--spatial-ensemble-max-order 2
--spatial-ensemble-min-ref-stars-per-frame 30
--spatial-ensemble-clip-sigma 3.5
--spatial-ensemble-n-iterations 3
--spatial-ensemble-min-bic-improvement 6
--spatial-ensemble-min-spatial-peak-to-peak 0.003
--spatial-ensemble-max-residual-mad 0.05
--spatial-ensemble-max-residual-mad-degradation 0.0
--spatial-ensemble-max-abs-spatial-deviation 0.15
--spatial-ensemble-min-good-frame-fraction 0.8
\end{lstlisting}
Параметр экспериментальный и по умолчанию выключен. Значение
\texttt{--spatial-ensemble-max-residual-mad-degradation 0.0} запрещает принимать spatial surface,
если её robust residual MAD хуже identity; отрицательное значение отключает только этот локальный
non-degradation gate.

\subsection*{Post-ensemble photometric frame QC и candidate axis correlation (0.5.44.16)}

Для консервативного контроля инструментальной систематики без изменения величин и ошибок SUSPECT-эпох:
\begin{lstlisting}
python -m grandgousier.cli process-archive-path ... \
  --photometry-estimator adaptive_ellipse \
  --enable-photometric-frame-qc
\end{lstlisting}
Параметр нельзя одновременно использовать с \texttt{--enable-spatial-ensemble}. Relative residual и
spatial anomalies дают SUSPECT/\texttt{KEEP\_FLAGGED}; только hard failures становятся BAD и
исключаются из standard science. Оценка excess frame scatter сохраняется как diagnostic-only и не
изменяет \texttt{mag\_err}, \texttt{mag\_err\_random} или \texttt{mag\_err\_systematic}.

Accepted automatic candidates дополнительно проверяются обычным candidate evaluator только на GOOD
кадрах. Используются четыре класса confirmation:
\begin{itemize}
  \item \texttt{GOOD\_FRAME\_CONFIRMED};
  \item \texttt{SYSTEMATICS\_SENSITIVE};
  \item \texttt{INSUFFICIENT\_GOOD\_FRAME\_SUPPORT};
  \item \texttt{TARGET\_PRESERVED}.
\end{itemize}
Диагностика сохраняется в \texttt{photometric\_frame\_qc/iterN} и в summary GOOD-frame
confirmation; PSF summary записывается в \texttt{frames.quality\_flags}.

Для \texttt{GOOD\_FRAME\_CONFIRMED} candidate final export не использует \texttt{HIGH\_RMS} как
point-level reject reason; сам флаг сохраняется в Object QC. Остальные reject flags не ослабляются.
Начиная с 0.5.44.15 candidate-triage CSV содержит GOOD-frame confirmation diagnostics, а PSF summary
сохраняет BIC-selected spatial model и peak-to-peak для geometric FWHM и ellipticity.

В 0.5.44.16 тот же режим автоматически добавляет diagnostic-only корреляцию accepted
\texttt{GOOD\_FRAME\_CONFIRMED} candidates внутри X/Y detector bins: 16 primary bins плюс
half-bin-shifted разметка и минимум 8 общих GOOD epochs. Начиная с 0.5.44.17 параллельно считаются
Spearman и Pearson; диагностический axis flag ставится при $|\rho_S|\ge0.8$ или $|r_P|\ge0.8$.
Обе signed-метрики и режим \texttt{SPEARMAN/PEARSON/BOTH/NONE} сохраняются в export.
Результат не меняет candidate status. Pair table: \texttt{candidate\_axis\_correlation/candidate\_axis\_pairs.csv}.


\section{P4a: stable physical-observation identity}

Read-only plan:

\begin{lstlisting}
python -m grandgousier.cli backfill-observation-identity `
  --instrument MASTER `
  --instrument RoboPhot `
  --out-root reports\observation_identity_055200
\end{lstlisting}

После проверки \texttt{collisions=0}, \texttt{invalid=0}, \texttt{link\_conflicts=0} и \texttt{can\_execute=YES} выполните ту же команду с \texttt{--execute}. MASTER identity v1 основана на \texttt{UNIQNUM}; RoboPhot identity v1 --- на \texttt{SERNUM + DATE-OBS}, где \texttt{DATE-OBS} уже является UTC. Backfill не объединяет и не удаляет строки \texttt{frames}.

\section{P4c: plan-only observation science-input guard}

Перед обработкой каталога, который может содержать копии уже известных observations, используйте обычный \texttt{process-archive-path} без \texttt{--execute} и сохраните manifest:

\begin{lstlisting}
python -m grandgousier.cli process-archive-path `
  --path "D:\archive\candidate_block" `
  --instrument-profile MASTER `
  --catalog-name "existing_shared_catalog" `
  --manifest-out "reports\p4c_plan.json"
\end{lstlisting}

Помимо JSON создаётся \texttt{reports\textbackslash p4c\_plan\_observation\_selection.csv}. Проверьте actions \texttt{ALREADY\_REPRESENTED}, \texttt{AMBIGUOUS\_MULTIPLE\_PRODUCTS}, \texttt{EXISTING\_MULTIPLE\_REPRESENTATIVES}, \texttt{MISSING\_OBSERVATION\_ID} и \texttt{IDENTITY\_CONFLICT} до запуска с \texttt{--execute}.

Начиная с 0.5.54.1 повторный обычный \texttt{process-archive-path --execute}, для которого guard не выбрал ни одного нового Frame, не перезапускает series, найденные только через уже представленные products. Для преднамеренного rerun существующей series используйте \texttt{--series-id}.


\section{Operational self-service: config, doctor, status}

Создать шаблон общего конфига:
\begin{lstlisting}
python -m grandgousier.cli config-example --out grandgousier.yaml
\end{lstlisting}

Проверить installation/deployment:
\begin{lstlisting}
python -m grandgousier.cli doctor --config grandgousier.yaml
python -m grandgousier.cli doctor --config grandgousier.yaml --strict-optional
\end{lstlisting}

Показать read-only состояние поля:
\begin{lstlisting}
python -m grandgousier.cli status `
  --config grandgousier.yaml `
  --field-id 33 `
  --catalog-name "auto_catalog_master_20180907_GPX_TF17_2_mixed" `
  --verbose
\end{lstlisting}

Опция \cli{--config} допустима и перед bootstrap subcommand: \cli{python -m grandgousier.cli --config grandgousier.yaml doctor}. Для \cli{process-field-archive} новый общий config использует вложенную секцию \texttt{field\_archive}; legacy \cli{--job-config} не удалён.


\section{Run-bundle diagnostics}

\begin{lstlisting}
grandgousier diagnose-run reports\field_archive\_jobs\<job_id>
grandgousier diagnose-run reports\field_archive\_jobs\<job_id>\job_summary.json --format json
\end{lstlisting}

Команда только читает локальные JSON/JSONL artifacts, не подключается к БД и не меняет processing state. Для bundle 0.5.59+ выводятся status, stable recovery code и рекомендуемые действия; legacy bundle также поддерживается.

\section{Scientific acceptance}

\begin{lstlisting}
grandgousier acceptance-check `
  --dataset gpx_tf17_2 `
  --config grandgousier.yaml `
  --report-out reports\acceptance_gpx_tf17_2.json
\end{lstlisting}

Команда read-only по PostgreSQL. Без \cli{--out-root} временные M5 products удаляются после сравнения; JSON report содержит все checks и фактически измеренные значения.
\subsection{Компактный научный экспорт Parquet (source candidate)}

Для Parquet требуется необязательная зависимость: \texttt{pip install grandgousier[parquet]}.
Продукт состоит из разреженной фотометрии SCIENCE, источников, кадров, сводки серий и
dataset-манифеста; отсутствующие source-frame точки не генерируются. Ночь наблюдений
определяется как дата UTC-времени \texttt{DATE-OBS} минус 12 часов. Это описание
контракта исходного кода, а не подтверждение успешного экспорта из рабочей БД.
